\section{Vulnerability Disclosure Policy：把外部研究变成修复输入}

外部研究人员可能比内部团队更早发现问题，但如果没有清晰报告渠道，他们不知道哪些系统可以测试、什么行为被允许、何时会收到回应。\term{Vulnerability Disclosure Policy}（VDP，漏洞披露政策）是一份公开规则与接收机制：定义 scope、授权边界、提交方式、组织承诺和协调沟通。它首先解决“如何安全报告”，并不保证支付奖励。

\subsection{从 report 到 coordinated disclosure}

\term{coordinated vulnerability disclosure}（协调漏洞披露，CVD）是研究人员、受影响组织和可能的供应商在修复与公开之间协调时间、影响范围和信息内容的过程。VDP 的工程价值在于把模糊邮件变成可追踪 workflow：acknowledge、triage、validate、scope、remediate、verify、communicate。这里的 service-level objective 不是为了漂亮 dashboard，而是降低研究人员重复测试、公开升级或失去信任的概率。

\lecturefigure{03-vdp-intake.png}{VDP 将外部研究报告转化为验证、修复与协调披露流程。}{CISA VDP template；本地字幕 12:30--15:20；概念重绘。}

\begin{knowledgebox}{读图：每个 handoff 都要返回状态}
Reporter 提交资产、影响和复现证据；triage 确认 scope、severity 与 duplicate；remediation 指定 owner、修复、验证和 timeline；close/disclose 向研究人员更新状态，并在需要时协调供应链或公众。最容易失败的是中间静默：组织已收到却不确认、已修复却不验证、需要延期却不给理由。透明不等于公开所有内部信息，而是建立可预测承诺。
\end{knowledgebox}

\begin{importantbox}{VDP intake 的最小字段}
\begin{itemize}
  \item 受影响 asset、endpoint 或 product version；
  \item 观察到的 security impact 与最小复现证据；
  \item 测试时间、researcher contact 和是否可能涉及第三方；
  \item 自动生成 case ID、acknowledgement 与状态更新窗口；
  \item severity、owner、containment、fix、verification 和 disclosure decision。
\end{itemize}
不要要求研究人员提交与验证无关的个人资料，也不要让 bounty payment 成为修复排期的唯一优先级。
\end{importantbox}

\teachervoice{Sullivan 把 responsible disclosure 描述为一项基础设施创新：它向善意研究者说明“如果你按这些边界报告，我们希望听到，并会通过既定渠道处理”。课堂提示是关系设计，而不只是 security.txt 地址。}

\begin{warningbox}{VDP 不能自动把事件降级为普通漏洞}
外部报告可能同时表明 unauthorized access、credential compromise、数据获取或持续威胁。此时 VDP case 仍负责 researcher communication，但事件必须进入独立 incident command、evidence preservation 与 disclosure analysis。把所有事实留在 bounty platform 会遮蔽严重性。
\end{warningbox}

\subsection{本章小结}

VDP 是安全报告渠道和信任合同。它通过 scope、响应承诺与协调修复降低双方不确定性，但一旦事实显示真实 compromise，就必须触发更强的事件治理。

